\begin{exbox}[连续模共振根的 Jacobian 与真空单光子检查]
设连续模按角频率归一化，$[a(\omega),a^\dagger(\omega')]=\delta(\omega-\omega')$。
若吸收项为 $\hbar g(\omega)B^\dagger a(\omega)\dd\omega$，则 $[g]={\rm s}^{-1/2}$。
取平稳且频率对角的光输入，满足
$\langle a^\dagger(\omega)a(\omega')\rangle=n(\omega)\delta(\omega-\omega')$，
其中 $n$ 无量纲。有限时间一阶吸收数谱为
\begin{equation}
 \frac{\dd P_+(T)}{\dd\omega}=|g(\omega)|^2 n(\omega)
 T^2\sinc^2[\Delta(\omega)T/2].
 \label{eq:feqo-continuum-mode-amplitude}
\end{equation}
左边单位为秒，和右边的 $|g|^2T^2$ 一致。
若共振根 $\omega_j$ 均为简单根，$\Delta'(\omega_j)\ne0$，长时谱求和给
\begin{equation}
 \Gamma_+=2\pi\sum_j\frac{|g(\omega_j)|^2n(\omega_j)}{|\Delta'(\omega_j)|}.
 \label{eq:feqo-continuum-absorption}
\end{equation}
证明是在每个根附近作变量替换 $u=\Delta(\omega)$，
$\dd\omega=\dd u/|\Delta'(\omega_j)|$，再对有限时间核使用近似恒等式。
无共振根时速率为零；多重根或带边处的 Jacobian 发散表示这个简单根黄金律失效，
应保留有限时间核和真实谱结构，而不能把发散速率直接当作实验概率。

还需同时施加动量约束。真空中自由光子 $K^\mu=(\omega/c,\mathbf k)$ 满足 $K^2=0$，
自由电子 $P^2=m_{\rm e}^2c^2$。若电子吸收一个真实光子，则末态为 $P+\hbar K$，
质量壳要求 $2\hbar P\cdot K=0$。
然而电子静止系中 $P\cdot K=m_{\rm e}\omega>0$，且此内积为 Lorentz 不变量，矛盾。
发射的末态 $P-\hbar K$ 也要求同一不可能条件。因此无外部结构、无其他粒子的自由电子，
不能仅凭吸收或发射一个真空真实光子构成守恒的渐近过程。
盒归一化中模式态密度的体积因子与矩阵元平方确会抵消，
但这种形式操作并不会创造一个原本不存在的质量壳交点。
近场、介质、周期结构或有物理开关的外源能提供额外动量与能量，
这才是前面有限窗口及结构谱能够产生真实信号的条件。
\end{exbox}
